\documentclass[b4paper,12pt]{exam}

%% 引入Package %%
\usepackage{graphicx}
\usepackage{extsizes}
\usepackage[margin=2cm]{geometry}
\usepackage{caption}
\usepackage{xcolor}
\usepackage{mathtools,amsthm,amssymb}
\usepackage[shortlabels,inline]{enumitem}
\usepackage{tikz}
\usepackage{pgfplots}
\usetikzlibrary{decorations.pathmorphing,decorations.markings,arrows.meta,calc,intersections,bending}
\usepackage{ulem}
\usepackage{enumitem}
\usepackage{ifthen}
\usepackage{draftwatermark}

% 浮水印開關
\newif\ifshowWatermark
\showWatermarktrue
\ifshowWatermark
  \SetWatermarkText{中山附中樣卷}
  \SetWatermarkScale{0.8}
\else
  \SetWatermarkText{}
\fi

% 選擇題排版指令
\newcommand{\anslabel}[1]{%
  \makebox[3em][c]{%
    \ifshowAnswer{\color{red}#1}\else\phantom{#1}\fi
  }%
}
\newcommand{\anslist}[1]{#1}

\newcommand{\fourch}[7]{\vspace{#6}\\\begin{tabular}{*{4}{@{}p{#7}}}(A)~#1 & (B)~#2 & (C)~#3 & (D)~#4 & (E)~#5\end{tabular}}
\newcommand{\twoch}[7]{\vspace{#6}\\\begin{tabular}{*{4}{@{}p{#7}}}(A)~#1 & (B)~#2\end{tabular}\vspace{#6}\\\begin{tabular}{*{4}{@{}p{#7}}}(C)~#3 & (D)~#4\end{tabular}\vspace{#6}\\\begin{tabular}{*{4}{@{}p{#7}}}(E)~#5 & ~\end{tabular}}
\newcommand{\onech}[6]{\vspace{#6}\\(A)~#1\vspace{#6}\\(B)~#2\vspace{#6}\\(C)~#3\vspace{#6}\\(D)~#4\vspace{#6}\\(E)~#5}

% 簡化命令
\newcommand\twoSolution[2]{$x=#1$ 或 $x=#2$}

% 標題區塊參數化
\newcommand\testTitle[1]{\section*{#1}}
\newcommand\testSubsection[4]{\subsection*{#1 (每#2\ #3\ 分，共\ #4\ 分)}}

% 中文字體 (需 XeLaTeX)
\usepackage[CJKmath=true,AutoFakeBold=3,AutoFakeSlant=.2]{xeCJK}
\setCJKmainfont{AR PL KaitiM Big5}
\setCJKsansfont{AR PL KaitiM Big5}
\setCJKmonofont{AR PL KaitiM Big5}
\newCJKfontfamily\Kai{標楷體}
\newCJKfontfamily\Hei{微軟正黑體}
\newCJKfontfamily\NewMing{新細明體}
\XeTeXlinebreaklocale "zh"

% 格式設定
\setlength{\headheight}{15pt}
\setlength{\parindent}{22pt}

% 答案開關
\newif\ifshowAnswer
\showAnswerfalse
\newcommand{\colorAnswer}[2]{%
{
  \makebox[6em][c]{% <-- 固定寬度，自己調整
    \ifshowAnswer
      {\color{#1}#2}%
    \else
      % 不顯示答案，但保持空間
      % 可以放空白，或放透明字
      \phantom{#2}%
    \fi
  }}%
}

\begin{document}
\captionsetup[figure]{labelfont={bf}, name={\text{圖}}, labelsep=period}
\captionsetup[table]{labelfont={bf}, name={\text{表}}, labelsep=period}

% ====== 考卷標題 ======
{\centering
\testTitle{國立中山大學附中114學年度第一學期高一上數學第四次周考}
}

\hfill 一年\hspace{2em}班\hspace{2em}號\quad 姓名：\hspace{6em}\quad\\
\vspace{0.5em}

\footer {}{\iflastpage{\textbf{請記得檢查並驗算}}{\oddeven{\textbf{背面仍有題目，請翻面作答}}{}}}{}

% ====== 一、多重選擇題 ======
\testSubsection{一、多重選擇題}{題}{8}{16}
\begin{enumerate}[itemsep=0.5em, topsep=1em,
                  label=\arabic*.]
    \item(\anslabel{CDE}) 下列敘述關於圓方程式哪些選項正確？\\[.5em]
        (A)~{已知$\Gamma_1:(x+h)^2-(y+k)^2=r^2$，則畫在座標平面上可得$\Gamma$為一個圓}\\[.5em]
        (B)~{已知圓$\Gamma_2:(x+h)^2+(y+k)^2=r^2$，則圓心在 $(h,k)$半徑為$r$。}\\[.5em]
        (C)~{圓方程式寫成$x^2+y^2-dx-ey+f=0$稱為圓的一般式。}\\[.5em]
        (D)~{若方程式$x^2+y^2-dx-ey+fxy+g=0$為圓方程式，則 $f=0$}\\[.5em]
        (E)~{方程式$\dfrac{x^2}{a}+\dfrac{y^2}{b}=1$ ，當 $a=b\neq0$ 時為圓方程式。}

    \item(\anslabel{ACE}) $\Gamma_3:ax^2+by^2+cxy+dx+ey+f=0$為圓方程式，且圓心在$(h,k)$，半徑為$r$，下列敘述何者正確？\\[.5em]
        (A)~{$a\neq0$}\hspace{2em}
        (B)~{$bc=1$}\hspace{2em}
        (C)~{$h=-\cfrac{d}{2a}$}\hspace{2em}
        (D)~{$r=\cfrac{\sqrt{d^2+e^2-f}}{2}$}\hspace{2em}
        (E)~{圓的判別式為 $\Delta=d^2+e^2-4f$}\\
\end{enumerate}

% ====== 二、填充計算題 ======
\subsection*{二、填充計算題 (每格\ 7\ 分，共\ 84\ 分)}
\textbf{* 若題目無特別註明，圓方程式請一律以標準式寫出，否則不予計分。}
\begin{enumerate}[itemsep=1em, topsep=2em]
    \item 已知圓心 $(14,20)$ 半徑為 $10$，試問圓方程式：\underline{\colorAnswer{red}{$?$}}。

    \item 寫出座標平面上與$(-3,1)$距離為$3\sqrt{2}$的所有點的方程式：\underline{\colorAnswer{red}{$?$}}。

    \item 已知圓 $\Gamma_4:x^2+y^2=9$，今將圖形向右移動 $5$ 單位，向上移動 $2$ 單位，試求新圓方程式：\underline{\colorAnswer{red}{$?$}}。
    
    \item 已知 $\triangle ABC$為圓$\Gamma_5$的內接三角形，其中 $\angle ABC=90^\circ$，且 $A(5,4)$、$B(11,-4)$，試求圓方程式：\underline{\colorAnswer{red}{$?$}}。

    \item 座標平面上一圓與$x$軸相切且圓心在$(1,9)$，試求圓方程式：\underline{\colorAnswer{red}{$?$}}。

    \item 若一圓通過三點 $A(0,0)$、$B(7,7)$、$C(0,6)$，則圓方程式為：\underline{\colorAnswer{red}{$?$}}。

    \item 若圓方程式 $\Gamma_6:(x-a)^2+(y-4)^2=136$且通過點 $(-16,-2)$，試問 $a=$\underline{\colorAnswer{red}{$?$}}。

    \item 若方程式 $x^2+y^2+10x-8y+k=0$ 沒有圖形，求實數$k$之範圍\underline{\colorAnswer{red}{$?$}}。

    \item 已知圓$C_1:x^2+y^2-4x-2y-18=0$的圓心$O_1$和圓$C_2:x^2+y^2+4x+4y-32=0$的圓心$O_2$，則此兩圓的連心線段長$\overline{O_1O_2}$為\underline{\colorAnswer{red}{$?$}}。

    \item 座標平面上，若方程式 $a(x^2+xy+x)+b(x^2-xy+x)-6x-4xy+6y^2+12y-24=0$的圖形為一圓，則此圓的半徑為\underline{\colorAnswer{red}{$?$}}。
    
    \vspace{1em}

    \begin{minipage}[t]{0.7\linewidth}
        \item 座標平面上，一圓與$y$軸相切於$(0,-8)$且通過點$(4,0)$，試求圓方程式：\underline{\colorAnswer{red}{$?$}}。 
        
        \item 如圖，今有一遊樂園試營運新開幕摩天輪，摩天輪搭乘最低點距離地面$5$公尺。今工作人員搭到地面$15$公尺的車廂$A$時看到位於正對面同高度的車廂$B$相距自己$60$公尺，當車廂$A$轉動到某高度時，又看到位於正對面同高度的$C$車廂也相距自己$60$公尺，試問此時車廂$A$距離地面高度\underline{\colorAnswer{red}{$?$}}公尺。
    \end{minipage}
    \hfill
    \begin{minipage}[t]{0.25\linewidth}
    \begin{tikzpicture}[baseline=(current bounding box.north), scale=0.5]
        \draw (0,0) circle (5);
        \draw[very thin] (0,0)--(-2,-5.5)--(2,-5.5)--(0,0);
        \draw[very thick] (-3.5,-3.5) coordinate[label=below left:$A$]--(3.5,-3.5) coordinate[label=below right:$B$];
        \draw[very thick] (-3.5,3.5) coordinate[label=above left:$A$]--(3.5,3.5) coordinate[label=above right:$C$];
        \begin{scope}[shift={(-1,0)}] % 只有這裡面的內容會被平移
            \draw[-latex, very thick] (-3.5,-3.5) arc (225:135:5);
        \end{scope}
    \end{tikzpicture}
    \end{minipage}

    % ...可繼續新增題目
\end{enumerate}

\end{document}
